\documentclass[12pt]{article}
\usepackage[utf8]{inputenc}
\usepackage[T5]{fontenc}
\usepackage[vietnamese]{babel}
\usepackage[a4paper,margin=2.5cm]{geometry}
\usepackage{setspace}
\usepackage{indentfirst}
\usepackage{titlesec}
\usepackage{hyperref}
\usepackage{graphicx}
\usepackage{amsmath,amssymb,amsthm,mathtools}
\usepackage{enumitem}

\setstretch{1.25}
\setlength{\parindent}{1cm}
\setlength{\parskip}{0.45em}

\titleformat{\section}
  {\normalfont\Large\bfseries}{\thesection}{1em}{}
\titlespacing{\section}{0pt}{1.2em}{0.7em}

\hypersetup{
    colorlinks=true,
    linkcolor=blue,
    filecolor=magenta,
    urlcolor=cyan,
    citecolor=blue
}

\title{Mô hình hóa bài toán motion planning \\ thành bài toán optimal control trên đồ thị các vùng lồi bằng phương pháp multiple shooting}
\author{}
\date{}

\begin{document}

\maketitle
\tableofcontents
\newpage

\section{Mục tiêu}

Mục tiêu là formulate một bài toán motion planning \textbf{tìm đồng thời}
\begin{enumerate}[label=(\roman*)]
    \item một \textbf{chuỗi rời rạc các vùng an toàn lồi}, và
    \item một \textbf{quỹ đạo liên tục theo thời gian} thỏa động lực học,
\end{enumerate}
sao cho toàn bộ quỹ đạo tránh vật cản và tối ưu một hàm chi phí toàn cục.

Phần \emph{discrete} (chọn vùng nào, theo thứ tự nào) và phần \emph{continuous} (quỹ đạo, điều khiển, thời lượng) phải được tối ưu \textbf{cùng một lúc}.

Về nguồn gốc ý tưởng, mô hình này ghép hai khối lý thuyết:
\begin{itemize}
    \item \textbf{Graphs of Convex Sets / network flow} để chọn một đường đi rời rạc qua các vùng an toàn lồi \cite{marcucci2021,marcucci2022};
    \item \textbf{direct multiple shooting} để mô hình hóa quỹ đạo liên tục trong mỗi vùng lồi.
\end{itemize}

Trong \emph{Shortest Paths in Graphs of Convex Sets}, quyết định tối ưu bao gồm đồng thời một đường đi rời rạc trên đồ thị và các biến liên tục gắn với những đỉnh/cạnh được kích hoạt \cite{marcucci2021}. Trong \emph{Motion Planning around Obstacles with Convex Optimization}, không gian collision-free được phân hoạch thành các safe convex regions, và mục tiêu là xây dựng một continuous trajectory nằm hoàn toàn trong hợp của các vùng ấy \cite{marcucci2022}. Còn trong \emph{direct multiple shooting}, state trên mỗi đoạn thời gian được thay bởi nghiệm của một local initial value problem bắt đầu từ một initial value, và các đoạn được nối lại bằng continuity defects \cite{bockplitt,diehl2005}.

\section{Bài toán gốc cần giải}

Xét một hệ động lực học liên tục
\begin{equation}
\dot x(t)=f(x(t),u(t)),
\qquad t\in[0,T],
\label{eq:ct_dyn}
\end{equation}
trong đó:
\begin{itemize}
    \item $x(t)\in\mathbb R^{n_x}$ là trạng thái,
    \item $u(t)\in\mathbb R^{n_u}$ là điều khiển,
    \item $T>0$ là thời điểm kết thúc (có thể cố định hoặc là biến quyết định).
\end{itemize}

Cho trạng thái đầu và trạng thái đích
\[
x(0)=x^{\mathrm s},
\qquad
x(T)=x^{\mathrm g},
\]
và cho không gian tự do (collision-free space)
\[
Q \subset \mathbb R^{n_x}.
\]

Trong toàn bộ báo cáo, để tránh nặng ký hiệu, ta viết trực tiếp $x(t)\in Q$. Nếu trạng thái $x$ còn chứa thêm các biến động lực không mang ý nghĩa hình học, thì điều kiện collision-free đúng cần được hiểu là $\pi(x(t))\in Q$ đối với một phép chiếu cấu hình $\pi$; các vùng $Q_v$ về sau cũng được hiểu theo cùng nghĩa đó.

Ta muốn giải bài toán optimal control tổng quát
\begin{align}
\min_{T,\,x(\cdot),\,u(\cdot)}\quad
& \Phi\bigl(x(T),T\bigr)+\int_0^T \ell\bigl(x(t),u(t)\bigr)\,dt
\label{eq:orig_ocp_obj}\\
\text{s.t.}\quad
& \dot x(t)=f(x(t),u(t)), && t\in[0,T],
\label{eq:orig_ocp_dyn}\\
& x(0)=x^{\mathrm s},\qquad x(T)=x^{\mathrm g},
\label{eq:orig_ocp_bc}\\
& x(t)\in Q,\qquad u(t)\in \mathcal U, && \forall t\in[0,T].
\label{eq:orig_ocp_path}
\end{align}


Trong hàm mục tiêu \eqref{eq:orig_ocp_obj}, 
$\Phi(x(T),T)$ là chi phí cuối (terminal cost), dùng để đánh giá trạng thái đích và/hoặc thời gian kết thúc; 
còn $\ell(x(t),u(t))$ là chi phí tức thời (running cost), được tích lũy trên toàn khoảng thời gian thông qua tích phân 
$\int_0^T \ell(x(t),u(t))\,dt$. 
Vì vậy, hàm mục tiêu gồm một phần đánh giá tại điểm cuối và một phần đánh giá toàn bộ quá trình chuyển động.

\section{Phân hoạch không gian tự do thành các vùng an toàn lồi}

Giả sử ta đã có một họ hữu hạn các vùng an toàn lồi
\[
\mathcal Q:=\{Q_v\}_{v\in V_r},
\qquad
Q_v\subseteq Q \subset \mathbb R^{n_x},
\]
trong đó mỗi $Q_v$ là một tập lồi đóng, bị chặn, không giao với vật cản.

Từ các vùng này, ta xây dựng một đồ thị có hướng
\[
G=(V,E),
\qquad
V=V_r\cup\{\sigma,\tau\},
\qquad
E=E_{rr}\cup E_s\cup E_t,
\]
trong đó:
\begin{itemize}
    \item $V_r$ là tập các đỉnh tương ứng với vùng lồi,
    \item $\sigma$ là đỉnh nguồn,
    \item $\tau$ là đỉnh đích.
\end{itemize}

Các cạnh giữa hai vùng lồi được tạo theo điều kiện adjacency hình học
\begin{equation}
E_{rr}:=\{(u,v)\in V_r\times V_r:\ Q_u\cap Q_v\neq \varnothing\}.
\label{eq:adjacency}
\end{equation}

Các cạnh biên được thêm riêng cho source và target:
\begin{equation}
E_s:=\{(\sigma,v):\ x^{\mathrm s}\in Q_v\},
\qquad
E_t:=\{(v,\tau):\ x^{\mathrm g}\in Q_v\}.
\label{eq:boundary_edges}
\end{equation}

Điều kiện \eqref{eq:adjacency} chỉ nói rằng hai vùng chạm nhau về mặt hình học, tức là tồn tại một trạng thái kết nối thuộc cả hai vùng. Nó \textbf{không} tự động suy ra dynamic feasibility khi còn có control bounds, phân bổ thời gian, ràng buộc nonholonomic, hoặc động lực học phi tuyến. Vì vậy, một đường đi rời rạc từ $\sigma$ đến $\tau$ chỉ nên được hiểu là một \textbf{skeleton hình học khả dĩ}; tính khả thi động lực học sẽ được quyết định bởi các ràng buộc continuous-time ở các mục sau.

\section{Biến quyết định của mô hình}

Ta chia biến quyết định thành ba lớp.

\subsection{Biến rời rạc: chọn đường đi trên đồ thị}

Với mỗi cạnh $e=(u,v)\in E$, đặt
\begin{equation}
y_{uv}\in\{0,1\}.
\label{eq:ydef}
\end{equation}
Ý nghĩa:
\begin{itemize}
    \item $y_{uv}=1$: cạnh $(u,v)$ được chọn, nghĩa là chuỗi vùng đi từ $u$ sang $v$,
    \item $y_{uv}=0$: cạnh đó không được dùng.
\end{itemize}

Với mỗi đỉnh vùng $v\in V_r$, đặt thêm biến kích hoạt đỉnh
\begin{equation}
p_v\in\{0,1\}.
\label{eq:pdef}
\end{equation}
Ý nghĩa:
\begin{itemize}
    \item $p_v=1$: vùng $Q_v$ thực sự được quỹ đạo đi qua,
    \item $p_v=0$: vùng $Q_v$ không xuất hiện trong nghiệm.
\end{itemize}


Biến $p_v$ về nguyên tắc có thể được suy ra từ các biến cạnh $y_{uv}$ thông qua
\[
p_v=\sum_{(u,v)\in E} y_{uv}=\sum_{(v,w)\in E} y_{vw}, \qquad v\in V_r.
\]
Tuy nhiên, việc giữ riêng biến kích hoạt đỉnh $p_v$ giúp diễn đạt gọn hơn các ràng buộc và chi phí gắn trực tiếp với vùng $Q_v$ và với local trajectory segment tương ứng.


\subsection{Biến continuous cục bộ trên mỗi vùng}

Với mỗi vùng $v\in V_r$, ta gắn một \textbf{local trajectory segment} trên miền chuẩn hóa $\tau\in[0,1]$. Các biến continuous cục bộ là:
\begin{align}
& s_v^- \in \mathbb R^{n_x} && \text{: entry shooting state của vùng } v, \nonumber\\
& s_v^+ \in \mathbb R^{n_x} && \text{: exit state của vùng } v, \nonumber\\
& \Delta_v \ge 0 && \text{: thời lượng local của đoạn trong vùng } v, \nonumber\\
& w_v \in \mathbb R^{m_v} && \text{: tham số điều khiển trên vùng } v. \nonumber
\end{align}

Ngoài ra, từ các biến này ta sinh ra quỹ đạo cục bộ
\[
x_v:[0,1]\to\mathbb R^{n_x},
\qquad
u_v:[0,1]\to\mathbb R^{n_u},
\]
qua một local initial value problem sẽ được định nghĩa ở phần sau.

\paragraph{Giải thích.}
Trong direct multiple shooting, ta không quyết định trực tiếp toàn bộ trạng thái trên toàn chân trời thời gian chỉ từ đầu mút, mà thay vào đó mỗi đoạn thời gian có một \emph{ initial value} riêng $s_v^-$.

\subsection{Biến liên kết giữa hai vùng liên tiếp}

Với mỗi cạnh nội bộ $(u,v)\in E_{rr}$, đặt một biến interface
\begin{equation}
z_{uv}\in\mathbb R^{n_x}.
\label{eq:zdef}
\end{equation}
Ý nghĩa:
\begin{itemize}
    \item nếu cạnh $(u,v)$ được chọn, thì $z_{uv}$ là trạng thái vật lý tại kết nối giữa hai đoạn cục bộ trong vùng $u$ và vùng $v$,
    \item do đó $z_{uv}$ phải nằm trong phần giao $Q_u\cap Q_v$.
\end{itemize}

\section{Ràng buộc network flow cho phần rời rạc}

\subsection{Ràng buộc nguồn và đích}

Ta yêu cầu đúng một cạnh đi ra khỏi nguồn và đúng một cạnh đi vào đích:
\begin{equation}
\sum_{(\sigma,v)\in E} y_{\sigma v}=1,
\qquad
\sum_{(v,\tau)\in E} y_{v\tau}=1.
\label{eq:source_sink_flow}
\end{equation}

Ý nghĩa: nghiệm phải bắt đầu ở một vùng đầu và kết thúc ở một vùng cuối.

\subsection{Bảo toàn luồng tại mỗi vùng}

Với mỗi $v\in V_r$, ta áp
\begin{equation}
\sum_{(u,v)\in E} y_{uv}
=
\sum_{(v,w)\in E} y_{vw}
=
p_v.
\label{eq:flow_balance}
\end{equation}

Ý nghĩa của từng thành phần trong \eqref{eq:flow_balance}:
\begin{itemize}
    \item $\sum_{(u,v)} y_{uv}=p_v$: nếu vùng $v$ được dùng thì phải có đúng một cạnh đi vào,
    \item $\sum_{(v,w)} y_{vw}=p_v$: nếu vùng $v$ được dùng thì phải có đúng một cạnh đi ra,
    \item nếu $p_v=0$ thì không cạnh nào vào/ra vùng $v$ được kích hoạt.
\end{itemize}

\subsection{Ràng buộc loại chu trình}

\begin{equation}
\sum_{e\in I_v^{\mathrm{out}}} y_e \le 1,
\qquad \forall v\in V_r,
\label{eq:acyclicity}
\end{equation}
trong đó $I_v^{\mathrm{out}}$ là tập các cạnh đi ra từ đỉnh $v$.

\section{Mô hình động lực học cục bộ theo direct multiple shooting}

\subsection{Tham số hóa điều khiển trên từng vùng}

Với mỗi vùng $v\in V_r$, điều khiển local được tham số hóa bởi một vector hữu hạn chiều $w_v$ dưới dạng
\begin{equation}
u_v(\tau)=\mathcal U_v(\tau;w_v),
\qquad \tau\in[0,1].
\label{eq:control_param}
\end{equation}

Ví dụ, $\mathcal U_v(\tau;w_v)$ có thể là:
\begin{itemize}
    \item điều khiển hằng trên đoạn,
    \item điều khiển piecewise constant,
    \item spline hoặc polynomial bậc thấp.
\end{itemize}

Ví dụ, giả sử $u_v(\tau)\in\mathbb R^2$ là vận tốc phẳng của robot trong vùng $v$. Nếu chọn điều khiển piecewise constant với hai đoạn con, ta có thể viết
\[
u_v(\tau)=
\begin{cases}
u_{v,1}, & 0\le \tau < \tfrac12,\\
u_{v,2}, & \tfrac12 \le \tau \le 1,
\end{cases}
\qquad
w_v=(u_{v,1},u_{v,2}).
\]
Khi đó, thay vì tối ưu trên toàn bộ một hàm $u_v(\cdot)$, ta chỉ cần tối ưu trên hai vector điều khiển $u_{v,1}$ và $u_{v,2}$.

Ý tưởng của \eqref{eq:control_param} là thay bài toán tối ưu trên không gian hàm $u_v(\cdot)$, vốn là một bài toán vô hạn chiều, bằng bài toán tối ưu trên một vector hữu hạn chiều $w_v$. Nói cách khác, thay vì xem $u_v$ là một hàm tùy ý trên $[0,1]$, ta chỉ cho phép nó thuộc một họ hàm đã chọn trước và được mô tả hoàn toàn bởi các tham số trong $w_v$. Khi đó, một khi $w_v$ được cố định thì toàn bộ tín hiệu điều khiển $u_v(\tau)$ trên đoạn thuộc vùng $v$ cũng được xác định.

Việc tham số hóa này có hai vai trò chính:
\begin{itemize}
    \item Nó giúp đưa bài toán optimal control về dạng nonlinear program hữu hạn chiều, là dạng phù hợp với direct multiple shooting và với các solver số thực tế.
    \item Nó cho phép điều chỉnh mức độ chi tiết của lớp điều khiển: nếu chọn điều khiển hằng trên đoạn thì mô hình gọn và tốn ít chi phí tính toán hơn, nhưng kém linh hoạt hơn; nếu chọn piecewise constant với nhiều đoạn con, hoặc spline hay đa thức bậc thấp, thì quỹ đạo biểu diễn được sẽ phong phú hơn, đổi lại số biến quyết định và độ khó tính toán cũng tăng lên.
\end{itemize}

\subsection{Local IVP trên từng vùng}

Sau khi chuẩn hóa thời gian về $[0,1]$, local trajectory trong vùng $v$ được xác định bởi bài toán giá trị đầu
\begin{equation}
\begin{cases}
\dfrac{d x_v}{d\tau}(\tau)=\Delta_v\,f\bigl(x_v(\tau),u_v(\tau)\bigr), & \tau\in[0,1],\\[4pt]
x_v(0)=s_v^-,\\
u_v(\tau)=\mathcal U_v(\tau;w_v).
\end{cases}
\label{eq:local_ivp}
\end{equation}

Từ nghiệm của \eqref{eq:local_ivp}, ta định nghĩa endpoint map
\begin{equation}
F_v(s_v^-,w_v,\Delta_v):=x_v(1).
\label{eq:Fv}
\end{equation}

Ánh xạ $F_v$ được gọi là endpoint map, vì nó ánh xạ dữ liệu đầu của một đoạn động lực học gồm trạng thái đầu $s_v^-$, tham số điều khiển $w_v$ và thời lượng $\Delta_v$ thành trạng thái cuối thu được sau khi giải bài toán giá trị đầu \eqref{eq:local_ivp} trên toàn đoạn chuẩn hóa $[0,1]$.

Khi đó, ràng buộc defect của multiple shooting là
\begin{equation}
s_v^+ - F_v(s_v^-,w_v,\Delta_v)=0,
\qquad \forall v\in V_r.
\label{eq:defect}
\end{equation}

\paragraph{Ý nghĩa.}
Ràng buộc \eqref{eq:defect} nói rằng exit state của vùng $v$ phải đúng bằng endpoint của local IVP bắt đầu từ entry shooting state $s_v^-$. Đây là continuity condition trong direct multiple shooting, đúng với công thức chuẩn kiểu
\[
s_{i+1}=x_i(t_{i+1};s_i,q_i)
\]

\section{Ràng buộc an toàn hình học: đoạn local phải nằm trong vùng lồi tương ứng}

Nếu vùng $v$ được kích hoạt, ta muốn local trajectory của nó nằm hoàn toàn trong vùng an toàn tương ứng:
\begin{equation}
x_v(\tau)\in Q_v,
\qquad \forall \tau\in[0,1],
\quad \text{khi } p_v=1.
\label{eq:region_constraint}
\end{equation}

Ý nghĩa của \eqref{eq:region_constraint} là:
\begin{itemize}
    \item ta không còn phải viết trực tiếp tránh obstacle ở mọi thời điểm,
    \item nhưng vì $Q_v$ là vùng an toàn, nếu mỗi local segment nằm trong vùng được gán cho nó thì toàn bộ quỹ đạo sẽ nằm trong hợp của các safe regions đã chọn,
    \item do đó quỹ đạo toàn cục collision-free theo continuous time.
\end{itemize}

Tuy nhiên, ở dạng hiện tại, \eqref{eq:region_constraint} mới chỉ là điều kiện continuous-time mong muốn, chứ chưa phải một dạng có thể đưa vào solver.

Đối với framework hiện tại, một hiện thực hữu hạn chiều tự nhiên của \eqref{eq:region_constraint} là dùng \textbf{sampling trên một lưới hữu hạn}. Cụ thể, chia miền chuẩn hóa $[0,1]$ của vùng $v$ thành $M_v$ khoảng con
\[
0=\tau_{v,0}<\tau_{v,1}<\cdots<\tau_{v,M_v}=1,
\]
và trước hết chỉ kiểm tra điều kiện an toàn tại các điểm lưới $\tau_{v,0},\dots,\tau_{v,M_v}$. Vì local trajectory $x_v(\cdot)$ đã được sinh bởi local IVP \eqref{eq:local_ivp} và bị ràng buộc bởi defect constraint \eqref{eq:defect}, ta có thể ép điều kiện an toàn bằng cách yêu cầu trạng thái tại các mesh points này đều nằm trong $Q_v$.

Nếu $Q_v$ là một polytope
\[
Q_v=\{x\in\mathbb R^{n_x}:A_vx\le b_v\},
\]
thì
\[
A_vx_v(\tau_{v,k})\le b_v-\varepsilon_v,
\qquad k=0,\dots,M_v,
\]
với $\varepsilon_v>0$ là safety margin.
\section{Ràng buộc ghép nối giữa phần discrete và phần continuous}

Phần này buộc network-flow path và local multiple-shooting segments khớp nhau.

Với mỗi cạnh nội bộ $(u,v)\in E_{rr}$, ta áp logic sau:
\begin{equation}
y_{uv}=1 \Longrightarrow
\begin{cases}
z_{uv}\in Q_u\cap Q_v,\\
s_u^+=z_{uv},\\
s_v^-=z_{uv}.
\end{cases}
\label{eq:edge_coupling}
\end{equation}

Ý nghĩa từng dòng trong \eqref{eq:edge_coupling}:
\begin{itemize}
    \item $z_{uv}\in Q_u\cap Q_v$: điểm giao phải là một trạng thái hình học hợp lệ cho cả hai vùng,
    \item $s_u^+=z_{uv}$: đoạn local trong vùng $u$ kết thúc tại điểm nối,
    \item $s_v^-=z_{uv}$: đoạn local trong vùng $v$ bắt đầu từ đúng điểm nối ấy.
\end{itemize}

Như vậy, nếu cạnh $(u,v)$ được chọn thì hai local IVPs ở $u$ và $v$ được kết nối lại với nhau tại cùng một trạng thái vật lý.

Trong mô hình hóa cuối, perspective chỉ xử lý gọn phần membership $z_{uv}\in Q_u\cap Q_v$; còn hai đẳng thức $s_u^+=z_{uv}$ và $s_v^-=z_{uv}$ vẫn phải được bật/tắt rõ ràng bằng một cơ chế indicator tương thích với $y_{uv}$.

\section{Ràng buộc biên với source và target}

Ta cần gắn local trajectory đầu tiên với trạng thái đầu, và local trajectory cuối cùng với trạng thái đích.

Với mọi cạnh $(\sigma,v)\in E_s$:
\begin{equation}
y_{\sigma v}=1 \Longrightarrow s_v^-=x^{\mathrm s}.
\label{eq:source_bc}
\end{equation}

Với mọi cạnh $(v,\tau)\in E_t$:
\begin{equation}
y_{v\tau}=1 \Longrightarrow s_v^+=x^{\mathrm g}.
\label{eq:target_bc}
\end{equation}

Ý nghĩa: vùng đầu tiên trong chuỗi phải nhận đúng initial condition, còn vùng cuối cùng phải kết thúc ở goal state.

\section{Hàm mục tiêu của mô hình}

Trong phần còn lại của báo cáo, ta chọn running cost theo dạng
\begin{equation}
\ell\bigl(x(t),u(t)\bigr)
:=
w_L\,\|\dot q(t)\|^2 + w_E\,\|u(t)\|^2,
\qquad
q(t):=\pi\bigl(x(t)\bigr),
\qquad
w_L,w_E>0.
\label{eq:running_cost_choice}
\end{equation}
Ở đây, $q(t)$ là phần cấu hình hình học của trạng thái; nếu $x(t)$ bản thân đã là biến cấu hình thì có thể hiểu đơn giản là $q(t)=x(t)$. Thành phần $w_L\,\|\dot q(t)\|^2$ phạt độ lớn vận tốc dịch chuyển, còn thành phần $w_E\,\|u(t)\|^2$ phạt năng lượng điều khiển.

Ta đặt local cost trên mỗi vùng $v$ là
\begin{equation}
\begin{aligned}
J_v(s_v^-,w_v,\Delta_v)
:={}& a\,\Delta_v \\
& + \int_0^1 \Delta_v\left(
w_L\,\left\|\frac{1}{\Delta_v}\frac{d q_v}{d\tau}(\tau)\right\|^2
+
w_E\,\|u_v(\tau)\|^2
\right)\,d\tau,
\end{aligned}
\label{eq:local_cost}
\end{equation}
trong đó $q_v(\tau):=\pi\bigl(x_v(\tau)\bigr)$ là phần cấu hình hình học của trạng thái local.

Khi đó objective toàn cục là
\begin{equation}
\min \sum_{v\in V_r} p_v\,J_v(s_v^-,w_v,\Delta_v).
\label{eq:global_cost}
\end{equation}

Ý nghĩa của từng thành phần:
\begin{itemize}
    \item $a\Delta_v$: phạt thời lượng spent trong vùng $v$,
    \item tích phân local: running cost của local segment,
    \item hệ số $p_v$: chỉ những vùng thực sự được kích hoạt mới đóng góp vào objective.
\end{itemize}


Hàm $J_v$ biểu diễn chi phí cục bộ của local trajectory segment trong vùng $v$. Thành phần $a\Delta_v$ dùng để phạt thời lượng của đoạn, còn thành phần tích phân biểu diễn running cost tích lũy dọc theo động lực học cục bộ. Do đã có hạng $a\Delta_v$, mô hình vẫn giữ được một áp lực trực tiếp lên tổng thời gian, trong khi running cost tập trung phạt độ dài chuyển động và độ lớn điều khiển.

Do đoạn thời gian thật của vùng $v$ là $[t_v^-,t_v^+]$ với độ dài $\Delta_v=t_v^+-t_v^-$, ta chuẩn hóa nó về miền $\tau\in[0,1]$ bằng phép đổi biến
\[
t=t_v^-+\Delta_v\tau.
\]
Vì thế
\[
dt=\frac{dt}{d\tau}d\tau=\Delta_v\,d\tau.
\]
Nói cách khác, hệ số $\Delta_v$ xuất hiện như Jacobian của phép đổi biến thời gian từ biến thật $t$ sang biến chuẩn hóa $\tau$. Do đó, một tích phân running cost theo thời gian thật trên đoạn $v$,
\[
\int_{t_v^-}^{t_v^+}\ell(x(t),u(t))\,dt,
\]
được viết lại dưới dạng
\[
\int_0^1 \Delta_v\left(
w_L\,\left\|\frac{1}{\Delta_v}\frac{d q_v}{d\tau}(\tau)\right\|^2
+
w_E\,\|u_v(\tau)\|^2
\right)\,d\tau.
\]

\section{Formulation thô ban đầu của bài toán ghép}

Gom các thành phần ý niệm lại, ta thu được formulation thô ban đầu sau:
\begin{align}
\min_{\substack{y,p,z,\\ \{s_v^-,s_v^+,w_v,\Delta_v\}_{v\in V_r}}}
\quad & \sum_{v\in V_r} p_v\,J_v(s_v^-,w_v,\Delta_v)
\label{eq:joint_problem_obj}\\
\text{s.t.}\quad
& \sum_{(\sigma,v)\in E} y_{\sigma v}=1,
\qquad
\sum_{(v,\tau)\in E} y_{v\tau}=1,
\label{eq:joint_problem_source_sink}\\
& \sum_{(u,v)\in E} y_{uv}
=
\sum_{(v,w)\in E} y_{vw}
=
p_v,
\qquad \forall v\in V_r,
\label{eq:joint_problem_flow}\\
& \sum_{e\in I_v^{\mathrm{out}}} y_e \le 1,
\qquad \forall v\in V_r,
\label{eq:joint_problem_acyc}\\
& x_v'(\tau)=\Delta_v f\bigl(x_v(\tau),u_v(\tau)\bigr),
\qquad \tau\in[0,1],
\quad \forall v\in V_r,
\label{eq:joint_problem_dyn}\\
& x_v(0)=s_v^-,
\qquad
u_v(\tau)=\mathcal U_v(\tau;w_v),
\quad \forall v\in V_r,
\label{eq:joint_problem_ivp}\\
& s_v^+ - F_v(s_v^-,w_v,\Delta_v)=0,
\qquad \forall v\in V_r,
\label{eq:joint_problem_defect}\\
& x_v(\tau)\in Q_v,
\qquad \forall \tau\in[0,1],\ \forall v\in V_r \text{ với } p_v=1,
\label{eq:joint_problem_region}\\
& y_{uv}=1 \Longrightarrow
\begin{cases}
z_{uv}\in Q_u\cap Q_v,\\
s_u^+=z_{uv},\\
s_v^-=z_{uv},
\end{cases}
\qquad \forall (u,v)\in E_{rr},
\label{eq:joint_problem_edge}\\
& y_{\sigma v}=1 \Longrightarrow s_v^-=x^{\mathrm s},
\qquad \forall (\sigma,v)\in E_s,
\label{eq:joint_problem_source}\\
& y_{v\tau}=1 \Longrightarrow s_v^+=x^{\mathrm g},
\qquad \forall (v,\tau)\in E_t,
\label{eq:joint_problem_target}\\
& y_{uv}\in\{0,1\},\qquad p_v\in\{0,1\},
\qquad \forall (u,v)\in E,\ \forall v\in V_r.
\label{eq:joint_problem_binary}
\end{align}


Trong \eqref{eq:joint_problem_obj}--\eqref{eq:joint_problem_binary}, ràng buộc \eqref{eq:acyclicity} được thêm vào để giới hạn số cạnh đi ra từ mỗi vùng.

\section{Perspective cho hình học lồi và động lực học}

\subsection{Vấn đề của formulation}
Nếu giữ nguyên bài toán \eqref{eq:joint_problem_obj}--\eqref{eq:joint_problem_binary} ở dạng trực tiếp, thì có ba vấn đề đại số dễ xuất hiện:
\begin{enumerate}
    \item objective kiểu
    \[
    \sum_{v\in V_r} p_v\,J_v(s_v^-,w_v,\Delta_v)
    \]
    tạo cảm giác ``bật/tắt'' local cost, nhưng về mặt mô hình hóa lại trộn một biến nhị phân với một hàm chi phí local nói chung không lồi;
    \item các biến continuous tại những vùng không được đi qua vẫn có thể tồn tại như các biến tự do nếu không có cơ chế deactivation rõ ràng;
    \item nếu viết ghép nối theo kiểu tích nhị phân--liên tục, chẳng hạn $y_{uv}z_{uv}$, thì sẽ sinh ra các bilinear constraints không cần thiết.
\end{enumerate}

Vì vậy, mô hình hóa được khuyến nghị là tách bài toán thành hai lớp:
\begin{enumerate}
    \item \textbf{lớp hình học lồi/on--off}: dùng \emph{homogenization/perspective} cho các tập lồi gắn với đỉnh và cạnh;
    \item \textbf{lớp động lực học multiple shooting}: giữ riêng dưới dạng các ràng buộc phi tuyến được bật/tắt bằng indicator/disjunction, không cố ép chúng vào một perspective form khi bản thân chúng không lồi.
\end{enumerate}

\subsection{Perspective cho các tập lồi}
Giả sử $C\subset\mathbb R^n$ là một tập lồi đóng. Ta định nghĩa homogenization của $C$ bởi
\begin{equation}
\widetilde C:=\{(x,\lambda):\lambda>0,\ x/\lambda\in C\},
\qquad
\operatorname{cl}(\widetilde C)
\text{ là closure của nó.}
\label{eq:homogenized-set}
\end{equation}

Ý nghĩa của \eqref{eq:homogenized-set} là:
\begin{itemize}
    \item nếu $\lambda=1$ thì ta quay về đúng $x\in C$,
    \item nếu $\lambda=0$ thì ta thu được trạng thái ``off'' theo đúng closure của homogenization,
\end{itemize}

Áp dụng vào bài toán hiện tại, với mỗi vùng $v\in V_r$, ta có thể viết các ràng buộc hình học on--off dưới dạng
\begin{equation}
(s_v^-,p_v)\in \operatorname{cl}(\widetilde{Q}_v),
\qquad
(s_v^+,p_v)\in \operatorname{cl}(\widetilde{Q}_v).
\label{eq:perspective-vertex-states}
\end{equation}

Tương tự, với mỗi cạnh nội bộ $(u,v)\in E_{rr}$, thay vì chỉ viết
\[
z_{uv}\in Q_u\cap Q_v
\quad\text{khi } y_{uv}=1,
\]
ta viết trực tiếp
\begin{equation}
(z_{uv},y_{uv})
\in
\operatorname{cl}\!\bigl(\widetilde{(Q_u\cap Q_v)}\bigr).
\label{eq:perspective-edge-interface}
\end{equation}
Khi đó:
\begin{itemize}
    \item nếu $y_{uv}=1$ thì \eqref{eq:perspective-edge-interface} khôi phục đúng điều kiện $z_{uv}\in Q_u\cap Q_v$;
    \item nếu $y_{uv}=0$ thì biến $z_{uv}$ bị ``triệt tiêu'' theo closure của homogenization và không còn là một biến tự do có ý nghĩa vật lý độc lập.
\end{itemize}

Perspective ở đây chỉ xử lý gọn phần \emph{membership} trên các tập lồi. Nó không tự động mã hóa các đẳng thức ghép nối $s_u^+=z_{uv}$ và $s_v^-=z_{uv}$; các equalities này vẫn phải được viết rõ dưới dạng on--off constraints gắn với $y_{uv}$.

\subsection{Bật/tắt local multiple-shooting block}
Với mỗi vùng $v$, local multiple-shooting block vẫn phải giữ nguyên dưới dạng
\begin{equation}
\begin{cases}
\dfrac{d x_v}{d\tau}(\tau)=\Delta_v\,f\bigl(x_v(\tau),u_v(\tau)\bigr), & \tau\in[0,1],\\[4pt]
x_v(0)=s_v^-,\\
u_v(\tau)=\mathcal U_v(\tau;w_v),\\[2pt]
s_v^+ - F_v(s_v^-,w_v,\Delta_v)=0.
\end{cases}
\label{eq:ms-block-repeat}
\end{equation}
Khối \eqref{eq:ms-block-repeat} nói chung là \textbf{phi lồi}, vì $F_v$ là endpoint map sinh ra bởi tích phân một ODE và phụ thuộc phi tuyến vào các biến quyết định. Paper \emph{Motion Planning around Obstacles with Convex Optimization} cũng nhấn mạnh rằng nếu thêm continuous-time dynamics vào mô hình thì các equality constraints tương ứng trở nên nonconvex, ngay cả với linear control systems khi shape và timing cùng là biến quyết định \cite{marcucci2022}.

Do đó, perspective \textbf{không} giải được phần này. Cách đúng là giữ nó như một \textbf{nonlinear dynamic block} và bật/tắt nó bằng indicator/disjunction.


Với mỗi vùng $v\in V_r$, ta dùng logic sau:
\begin{equation}
p_v=1
\Longrightarrow
\begin{cases}
\Delta_v \ge \underline \Delta_v >0,\\
\text{local IVP } \eqref{eq:local_ivp} \text{ được kích hoạt},\\
s_v^+ - F_v(s_v^-,w_v,\Delta_v)=0,\\
\mathcal S_v^{\mathrm{safe}}(s_v^-,w_v,\Delta_v)\le 0,\\
\rho_v \ge \widehat J_v(s_v^-,w_v,\Delta_v),
\end{cases}
\label{eq:active-vertex-block}
\end{equation}
trong đó $\rho_v$ là epigraph variable cho local dynamic cost trên vùng $v$, $\widehat J_v$ là một quadrature approximation của chi phí đó, còn $\mathcal S_v^{\mathrm{safe}}$ trong báo cáo này được hiểu trước hết như một sampling-based safety certificate trên một lưới hữu hạn, có thể được siết dần bằng adaptive refinement để xấp xỉ \eqref{eq:region_constraint}.

Ngược lại, khi vùng $v$ không được dùng, ta triệt tiêu toàn bộ block:
\begin{equation}
p_v=0
\Longrightarrow
\begin{cases}
\Delta_v = 0,\\
\rho_v = 0,\\
s_v^- = 0,\\
s_v^+ = 0,\\
w_v = \bar w_v,
\end{cases}
\label{eq:inactive-vertex-block}
\end{equation}
trong đó $s_v^\pm=0$ là lựa chọn off-state tương thích với \eqref{eq:perspective-vertex-states}, còn $\bar w_v$ là một giá trị neo cố định đã chọn trước, thường lấy bằng $0$.

Ý nghĩa của \eqref{eq:active-vertex-block}--\eqref{eq:inactive-vertex-block} là:
\begin{itemize}
    \item nếu vùng active thì multiple-shooting dynamics và local cost thật sự được áp lên;
    \item nếu vùng inactive thì mọi biến local bị khóa vào trạng thái off và không thể ``trôi tự do'' trong mô hình.
\end{itemize}

\subsection{Mô hình cuối}
Không nên xem
\[
\sum_v p_v J_v(s_v^-,w_v,\Delta_v)
\]
là formulation solver-ready cuối cùng. Sau khi chỉnh lại các điểm ở trên, framework khuyến nghị gồm bốn lớp:
\begin{enumerate}
    \item \textbf{lớp discrete path}:
    \begin{align*}
    & \eqref{eq:source_sink_flow},\ \eqref{eq:flow_balance},\ \eqref{eq:acyclicity};
    \end{align*}
    \item \textbf{lớp convex/on--off geometry}:
    \begin{align*}
    & (s_v^-,p_v),(s_v^+,p_v) \in \operatorname{cl}(\widetilde{Q}_v),\qquad \forall v\in V_r,\\
    & (z_{uv},y_{uv})\in \operatorname{cl}\!\bigl(\widetilde{(Q_u\cap Q_v)}\bigr),\qquad \forall (u,v)\in E_{rr};
    \end{align*}
    \item \textbf{lớp ghép nối on--off}:
    \begin{align*}
    & y_{uv}=1 \Longrightarrow s_u^+=z_{uv},\qquad y_{uv}=1 \Longrightarrow s_v^-=z_{uv},
    \qquad \forall (u,v)\in E_{rr},\\
    & y_{\sigma v}=1 \Longrightarrow s_v^-=x^{\mathrm s},\qquad \forall (\sigma,v)\in E_s,\\
    & y_{v\tau}=1 \Longrightarrow s_v^+=x^{\mathrm g},\qquad \forall (v,\tau)\in E_t;
    \end{align*}
    \item \textbf{lớp nonlinear multiple-shooting dynamics}:
    \begin{align*}
    & p_v=1 \Longrightarrow
    \begin{cases}
    \text{local IVP } \eqref{eq:local_ivp},\\
    s_v^+ - F_v(s_v^-,w_v,\Delta_v)=0,\\
    \mathcal S_v^{\mathrm{safe}}(s_v^-,w_v,\Delta_v)\le 0,\\
    \rho_v \ge \widehat J_v(s_v^-,w_v,\Delta_v),
    \end{cases}
    \\
    & p_v=0 \Longrightarrow
    \begin{cases}
    \Delta_v=0,\ \rho_v=0,\\
    s_v^- = s_v^+ = 0,\\
    w_v=\bar w_v.
    \end{cases}
    \end{align*}
\end{enumerate}
Khi đó objective toàn cục được viết gọn hơn thành
\begin{equation}
\min \sum_{v\in V_r} \rho_v,
\label{eq:recommended-clean-objective}
\end{equation}
trong đó:
\begin{itemize}
    \item $\rho_v$ chứa local dynamic cost của multiple-shooting segment trên vùng $v$;
    \item block chi phí này được bật/tắt bằng indicator/disjunction đồng bộ với activation variable $p_v$.
\end{itemize}

\paragraph{Kết luận}
Perspective function được dùng để xử lý phần on--off convex của mô hình, đặc biệt là:
\begin{itemize}
    \item triệt tiêu các biến continuous tại các vùng/cạnh không active,
    \item viết gọn các interface constraints trên các tập lồi giao nhau.
\end{itemize}
Ở đây, ``interface constraints'' là các ràng buộc ghép nối tại điểm chuyển tiếp giữa hai vùng liên tiếp, chẳng hạn điều kiện $z_{uv}\in Q_u\cap Q_v$ cùng với $s_u^+=z_{uv}$ và $s_v^-=z_{uv}$. Dùng perspective cho biến kết nối giúp biểu diễn gọn logic ``cạnh active thì điểm nối phải nằm trong phần giao, cạnh inactive thì biến kết nối bị triệt tiêu'', thay vì phải viết nhiều ràng buộc on--off rời rạc hoặc sinh thêm các tích nhị phân--liên tục không cần thiết.
Tuy nhiên, perspective \textbf{không} loại bỏ được bản chất phi lồi của multiple shooting, cũng không thay thế cho cycle-elimination hay safety transcription hữu hạn chiều. Vì vậy, mô hình sau khi chỉnh đúng hơn nên được xem là một \textbf{mixed-integer nonlinear optimal-control framework} ở mức continuous time.

\section{Phân loại bài toán}

Mô hình trên \textbf{không còn là MISOCP} như GCS-B\'ezier. Lý do là các ánh xạ
\[
F_v(s_v^-,w_v,\Delta_v)
\]
được tạo ra bởi việc tích phân một ODE phi tuyến nói chung, và các ràng buộc vùng liên tục theo thời gian cũng không còn là thuần convex algebraic constraints.

Do đó, cách gọi chính xác hơn là:
\begin{equation*}
\boxed{
\begin{array}{c}
\text{Continuous-time level: Mixed-Integer Nonlinear Optimal Control Problem (MIOCP),}\\
\text{after finite-dimensional transcription: Mixed-Integer Nonlinear Program (MINLP).}
\end{array}}
\end{equation*}

\begin{thebibliography}{99}

\bibitem{bockplitt}
H.~G. Bock and K.~J. Plitt,
\newblock \emph{A Multiple Shooting Algorithm for Direct Solution of Optimal Control Problems},
\newblock Proceedings of the 9th IFAC World Congress, Budapest, 1984.

\bibitem{diehl2005}
M. Diehl, H.~G. Bock, H. Diedam, and P.-B. Wieber,
\newblock \emph{Fast Direct Multiple Shooting Algorithms for Optimal Robot Control},
\newblock in \emph{Fast Motions in Biomechanics and Robotics}, Springer, 2006.

\bibitem{marcucci2021}
T. Marcucci, J. Umenberger, P.~A. Parrilo, and R. Tedrake,
\newblock \emph{Shortest Paths in Graphs of Convex Sets},
\newblock SIAM Journal on Optimization, 2023.

\bibitem{marcucci2022}
T. Marcucci, M. Petersen, D. von Wrangel, and R. Tedrake,
\newblock \emph{Motion Planning around Obstacles with Convex Optimization},
\newblock The International Journal of Robotics Research, 2023.

\bibitem{marcucciThesis}
T. Marcucci,
\newblock \emph{Graphs of Convex Sets with Applications to Optimal Control and Motion Planning},
\newblock PhD thesis, Massachusetts Institute of Technology, 2024.

\end{thebibliography}

\end{document}
